\subsection{Measure induced on a locally compact subspace}

Let X be a locally compact space, \(\mu\) a measure on X, and Y a locally compact subspace of X. Since Y is the intersection of an open set and a closed set in X (GT, I, §9, No.~7, Prop. 12), it is \(\mu\) -measurable (No.~1, Cor. of Prop. 3). For every function \(g \in \mathcal{K}(\mathrm{Y};\mathbf{C})\) let \(g'\) be the function, defined on all of X, equal to g on Y and to 0 on X - Y; let us show that \(g'\) is \(\mu\) -integrable. We can restrict ourselves to the case that g is real and \(\geqslant 0\) (on writing g as a linear combination of such functions); since \(g'\) is bounded and has compact support, it suffices to show that \(g'\) is \(\mu\) -measurable (No.~6, Th. 5); but this follows from the fact that \(g'\) is upper semi-continuous on X (No.~5, Cor. of Prop. 8). We may therefore make the following definition:

DEFINITION 4. --- Given a locally compact subspace Y of a locally compact space X, we call measure induced on Y by a measure \(\mu\) on X, and denote by \(\mu_{Y}\) or \(\mu|Y\) , the measure defined by the formula

\[
\int g d \mu_ {\mathrm{Y}} = \int g ^ {\prime} d \mu\tag{1}
\]

for every function \(g \in \mathcal{K}(\mathrm{Y};\mathbf{C})\) , where \(g'\) denotes the function equal to \(g\) on \(\mathrm{Y}\) and to 0 on \(\mathrm{X} - \mathrm{Y}\) .

Example. --- Let \(\mu\) be Lebesgue measure on R, I any interval in R; I is a locally compact subspace of R, and the measure induced by \(\mu\) on I is the linear form

\[
g \mapsto \int_ {a} ^ {b} g (x) d x
\]

on \(\mathcal{K}(\mathrm{I};\mathbf{C})\) , where \(a\) and \(b\) are the endpoints (finite or not) of I (cf.~§ 4, No.~4, Example), in other words, what we have called Lebesgue measure on I.

When Y is an open subspace of X, Def. 4 coincides with the definition of the measure induced by \(\mu\) on Y (or the restriction of \(\mu\) to Y) given in Ch. III, §2, No.~1: indeed, for every function \(g \in \mathcal{K}(\mathrm{Y}; \mathbf{C})\) the function \(g'\) is then continuous on X.

We shall study integration with respect to an induced measure in detail in Ch. V, §7, and until then we shall need only the following results:

Lemma 2. --- Let \(\mu\) be a positive measure on \(X\) , and let \(K\) be a compact subset of \(X\) .

\begin{enumerate}
\def\labelenumi{(\roman{enumi})}
\item
  For every compact (resp. open) subset H of K, \(\mu_{\mathrm{K}}(\mathrm{H}) = \mu (\mathrm{H})\)
\item
  For a subset \(\mathbf{N}\) of \(\mathbf{K}\) to be \(\mu_{\mathbf{K}}\) -negligible, it is necessary and sufficient that it be \(\mu\) -negligible.
\item
  If S is the support of \(\mu_{\mathrm{K}}\) , then \(\operatorname{Supp}(\mu_{\mathrm{S}}) = \mathrm{S}\) .
\setcounter{enumi}{0}
\item
  We can restrict ourselves to the case that H is compact. Denote by f the characteristic function of H in the space K; f is upper semi-continuous, hence is the lower envelope of a decreasing directed family \((g_{\alpha})\) of functions in \(\mathcal{K}_{+}(\mathrm{K})\) ; we have \(\mu_{\mathrm{K}}(\mathrm{H}) = \inf_{\alpha} \int g_{\alpha} d\mu_{\mathrm{K}}\) (§4, No.~4, Cor. 2 of Prop. 5). If \(g_{\alpha}^{\prime}\) is the function equal to \(g_{\alpha}\) on K and to 0 on X - K, then \(g_{\alpha}^{\prime}\) is upper semi-continuous, and the lower envelope of the decreasing directed family \((g_{\alpha}^{\prime})\) is the characteristic function \(\varphi_{H}\) of H in the space X; therefore
\end{enumerate}

\[
\mu (\mathrm{H}) = \inf _ {\alpha} \int g _ {\alpha} ^ {\prime} d \mu = \inf _ {\alpha} \int g _ {\alpha} d \mu_ {\mathrm{K}} = \mu_ {\mathrm{K}} (\mathrm{H})
\]

by (1).

\begin{enumerate}
\def\labelenumi{(\roman{enumi})}
\setcounter{enumi}{1}
\item
  If N is \(\mu\) -negligible then, for every \(\varepsilon > 0\) , there exists a relatively compact open neighborhood U of N in X such that \(\mu(\mathrm{U}) \leqslant \varepsilon\) ; since \(\mathrm{K} - (\mathrm{U} \cap \mathrm{K})\) is compact, it follows from (i) that \(\mu_{\mathrm{K}}(\mathrm{U} \cap \mathrm{K}) \leqslant \mu(\mathrm{U}) \leqslant \varepsilon\) , thus N is \(\mu_{K}\) -negligible. Conversely, if N is \(\mu_{K}\) -negligible then there exists an open neighborhood V of N in X such that \(\mu_{\mathrm{K}}(\mathrm{V} \cap \mathrm{K}) \leqslant \varepsilon\) ; by (i), we have \(\mu(\mathrm{V} \cap \mathrm{K}) = \mu_{\mathrm{K}}(\mathrm{V} \cap \mathrm{K})\) , therefore \(\mu(\mathrm{N}) = 0\) since \(\varepsilon\) is arbitrary.
\item
  For every open set U in K that intersects S, by hypothesis \(\mu_{\mathrm{K}}(\mathrm{U}\cap\mathrm{S})\neq0\) , therefore \(\mu(\mathrm{U}\cap\mathrm{S})\neq0\) by (i), and since \(U\cap S\) is open in S, \(\mu_{\mathrm{S}}(\mathrm{U}\cap\mathrm{S})\neq0\) by (i); since every nonempty open set in S is of the form \(U\cap S\) , where U is open in K, this proves that \(\operatorname{Supp}(\mu_{\mathrm{S}})=\mathrm{S}\) .
\end{enumerate}

Lemma 3. --- Let Y be a locally compact subspace of X; for every measure \(\mu\) on X, \(|\mu_{\mathrm{Y}}| = |\mu|_{\mathrm{Y}}\) .

Let \(f\) be a function in \(\mathcal{K}_{+}(\mathrm{Y})\) and \(\varepsilon\) any number \(>0\) ; by definition, there exists a function \(g \in \mathcal{K}(\mathrm{Y};\mathbf{C})\) such that \(|g| \leqslant f\) and \(|\mu_{\mathrm{Y}}|(f) \leqslant |\mu_{\mathrm{Y}}(g)| + \varepsilon\) . Denoting by \(f'\) and \(g'\) the functions obtained by extending \(f\) and \(g\) , respectively, to be 0 on \(\mathrm{X} - \mathrm{Y}\) , we have \(\mu_{\mathrm{Y}}(g) = \mu(g')\) and, since \(|g'| \leqslant f'\) ,

\[
\left| \mu \left(g ^ {\prime}\right) \right| \leqslant | \mu | \left(\left| g ^ {\prime} \right|\right) \leqslant | \mu | \left(f ^ {\prime}\right) = | \mu | _ {Y} (f),
\]

whence \(|\mu_{\mathrm{Y}}|(f)\leqslant |\mu |_{\mathrm{Y}}(f) + \varepsilon\) and, since \(\varepsilon\) is arbitrary,

\[
| \mu_ {\mathrm{Y}} | (f) \leqslant | \mu | _ {\mathrm{Y}} (f).
\]

On the other hand, let K be the support of f and let U be a compact neighborhood of K in X such that \(|\mu|(\mathrm{U}-\mathrm{K}) \leqslant \varepsilon\) ; by Urysohn's theorem, there exists a function \(f_{1} \in \mathcal{K}_{+}(\mathrm{X})\) , extending f, with support contained in U and such that \(\|f_{1}\| = \|f\|\) . There exists a function \(h_{1} \in \mathcal{K}(\mathrm{X}; \mathbf{C})\) such that \(|h_{1}| \leqslant f_{1}\) and \(|\mu|(f_{1}) \leqslant |\mu(h_{1})| + \varepsilon\) . If h is the restriction of \(h_{1}\) to Y, then \(h \in \mathcal{K}(\mathrm{Y}; \mathbf{C})\) , \(|h| \leqslant f\) and \(\mu(h_{1}) - \mu_{\mathrm{Y}}(h) = \mu(h_{1}\varphi_{\mathrm{U}-\mathrm{K}})\) , therefore

\[
\left| \mu \left(h _ {1}\right) - \mu_ {\mathrm{Y}} (h) \right| \leqslant \| f \| \cdot | \mu | (\mathrm{U} - \mathrm{K}) \leqslant \varepsilon \| f \|;
\]

moreover, we have similarly

\[
| \mu | (f _ {1}) - | \mu | _ {\mathrm{Y}} (f) = | \mu | (f _ {1} \varphi_ {\mathrm{U} - \mathrm{K}}) \quad \text { and } \quad \left| | \mu | (f _ {1}) - | \mu | _ {\mathrm{Y}} (f) \right| \leqslant \varepsilon \| f \|.
\]

From this we infer that

\[
| \mu | _ {\mathrm{Y}} (f) \leqslant | \mu_ {\mathrm{Y}} (h) | + \varepsilon (1 + 2 \| f \|) \leqslant | \mu_ {\mathrm{Y}} | (f) + \varepsilon (1 + 2 \| f \|)
\]

and since \(\varepsilon\) is arbitrary, \(|\mu|_{\mathrm{Y}}(f) \leqslant |\mu_{\mathrm{Y}}|(f)\) , which completes the proof.

